\documentclass[12pt]{article}
\usepackage[utf8]{inputenc}
\usepackage[T1]{fontenc}
\usepackage[margin=1in]{geometry}
\usepackage{setspace}
\usepackage{xcolor}
\usepackage[colorlinks=true, urlcolor=blue!60!black, linkcolor=blue!60!black]{hyperref}
\setlength{\parindent}{2em}
\setlength{\parskip}{1.25em}
\renewcommand{\baselinestretch}{1.0}
\begin{document}
\begin{center}
{\LARGE \textit{United States Steel Corp. v. Multistate Tax Commission}, 434 U.S. 452 (1978)\par}
\vspace{0.8em}
{\large Case Brief --- Four-Page Edition\par}
\end{center}
\vspace{1em}

\section{Memory Jogger}

Applying the functional test of \textit{Virginia v.\ Tennessee}, the Court held 7--2 that the Multistate Tax Compact needed no congressional consent because it gave member states no power they lacked individually and did not enhance state power at the expense of federal supremacy (\href{https://supreme.justia.com/cases/federal/us/434/452/}{Justia}).

\section{Detailed Case Facts}

After \textit{Northwestern States Portland Cement Co.\ v.\ Minnesota} (1959) allowed states to tax apportioned income from interstate operations, Congress enacted nexus rules and commissioned the Willis Committee study. When federal bills stalled, state tax administrators drafted the Multistate Tax Compact in 1966 to promote uniformity, improve compliance, and avoid duplicative taxation. It took effect in 1967 upon adoption by seven states; 21 had joined by 1972. Twelve bills seeking congressional consent failed.

The Compact created the Multistate Tax Commission. Its uniform regulations were strictly advisory, and each state could reject them. Under optional Article VIII, a state could ask the Commission to audit corporations and seek compulsory process in other adopting states' courts, while each state kept sole power over rates, bases, assessments, and collection. Any state could withdraw by repealing its statute.

\section{Procedural History}

In 1972 United States Steel and other corporations filed a class action in the Southern District of New York on behalf of multistate taxpayers threatened with audits, alleging violations of the Compact Clause, the Commerce Clause, and the Fourth and Fourteenth Amendments. A three-judge district court granted summary judgment for the Commission. Abandoning the Fourth Amendment claim, the plaintiffs appealed directly; the Court affirmed on February 21, 1978.

\section{Judicial Votes}

Affirmed, \textbf{7--2}.
\begin{itemize}
\item \textbf{Majority:} Justice Powell, joined by Chief Justice Burger and Justices Brennan, Stewart, Marshall, Rehnquist, and Stevens.
\item \textbf{Dissenting:} Justice White, joined by Justice Blackmun.
\end{itemize}

\section{Holding}

The Multistate Tax Compact does not require congressional consent under the Compact Clause, U.S.\ Const.\ art.\ I, \S\ 10, cl.\ 3. Reaffirming \textit{Virginia v.\ Tennessee} (1893), consent is required only for agreements directed to forming a combination tending to increase state political power in a way that may encroach upon the just supremacy of the United States. The Commerce Clause and Fourteenth Amendment challenges also failed.

\section{Analysis of Opinions}

\textbf{Majority (Powell).} A literal reading would require consent for every trivial agreement. The Framers' distinctions among ``treaty,'' ``compact,'' and ``agreement'' had been lost, and doctrine settled on Justice Field's functional test, applied as a holding in \textit{New Hampshire v.\ Maine} (1976). The inquiry concerns \textit{potential} impact on federal supremacy and applies to formal compacts and reciprocal legislation alike, but neither multilateral membership nor an administrative body triggers it.

The Compact failed to trigger consent for three reasons. First, whatever leverage states gained against corporations, the test is power relative to the Federal Government: the Compact authorized nothing states could not do alone, delegated no sovereign power, and allowed withdrawal. Second, unitary-business audits, apportionment definitions, and reciprocal subpoenas were all methods each state could adopt itself, and Congress retained plenary power to preempt. Third, pressure on nonmember states reflected lawful fiscal competition. Footnote 33 added that a ``federal interest'' or enhanced lobbying clout does not threaten supremacy while Congress's power remains plenary. The harassment claims were unpreserved and irrelevant because only states could assess taxes.

\textbf{White, with Blackmun (dissenting).} The majority paid lip service to ``potential'' encroachment while demanding proof of actual harm. Congress had refused twelve consent bills, and the Commission lobbied to ``bottle up'' federal bills and opposed a U.S.--U.K.\ tax treaty. Commission audits, apportionment options, the reciprocal subpoena bloc, and the seven-state trigger were synergistic powers, and the Clause also protects nonmember states. Upholding so pervasive an institution left ``very little life remaining'' in the Clause.

\textbf{Later developments.} \textit{Northeast Bancorp v.\ Board of Governors} (1985) applied \textit{U.S.\ Steel} to uphold regional banking statutes and refined the indicia of a ``compact.''

\section{Examples: Future Applications}

\begin{enumerate}
\item \textbf{Same-Side: Advisory AI Commission.} Fifteen states create a commission that drafts model AI rules and conducts voluntary audits; states may reject rules and withdraw. No sovereign delegation, so no consent is needed.
\item \textbf{Same-Side: Teaching Licenses.} Twenty states recognize one another's licenses with an advisory council. Each could do so alone, so no consent is needed.
\item \textbf{Opposite-Side: Atlantic Maritime Compact.} Twelve states create a port authority that levies tariffs on foreign vessels, negotiates with European nations, and bars withdrawal for 25 years. It encroaches on federal foreign-affairs supremacy and requires consent.
\item \textbf{Opposite-Side: Transferring a County.} Two states move a populous county, shifting House seats and electoral votes. That alters national representation and requires consent.
\item \textbf{Fence-Sitter: National Popular Vote Compact.} Each state controls its electors and may withdraw, but a bloc controlling presidential elections may change the national structure and dilute nonmembers' influence, invoking White's view.
\end{enumerate}

\section{Critique}

\textbf{Nullification of text.} Originalist and textualist critics echo White: reading ``any Agreement or Compact'' so narrowly means virtually no domestic interstate arrangement requires consent unless it already violates another constitutional limit.

\textbf{Horizontal federalism and synergy.} Testing power only \textit{quoad} the Federal Government neglects the danger that a regional cartel of states imposes externalities on excluded sister states. And treating a permanent, staffed agency as identical to 21 states acting independently ignores collective-action dynamics and coordinated lobbying power.

\section{Key Quotations}

\begin{quote}
``But the test is whether the Compact enhances state power \textit{quoad} the National Government. This pact does not purport to authorize the member States to exercise any powers they could not exercise in its absence. Nor is there any delegation of sovereign power to the Commission; each State retains complete freedom to adopt or reject the rules and regulations of the Commission.''
\end{quote}
\noindent --- Justice Powell, Opinion of the Court, 434 U.S. at 473.

\begin{quote}
``Read literally, the Compact Clause would require the States to obtain congressional approval before entering into any agreement among themselves, irrespective of form, subject, duration, or interest to the United States.''
\end{quote}
\noindent --- Justice Powell, Opinion of the Court, 434 U.S. at 459.

\begin{quote}
``The Compact Clause, however, is directed to joint action by more than one State. If its only purpose in the present context were to require the consent of Congress to agreements between States that would otherwise violate the Commerce Clause, it would have no independent meaning. The Clause must mean that some actions which would be permissible for individual States to undertake are not permissible for a group of States to agree to undertake.''
\end{quote}
\noindent --- Justice White, dissenting, 434 U.S. at 482.

\end{document}
